\documentclass[11pt]{article}
\usepackage[top=1.00in, bottom=1.0in, left=1in, right=1in]{geometry}
\renewcommand{\baselinestretch}{1.1}
\usepackage{graphicx}
\usepackage{natbib}
\usepackage{amsmath}
\usepackage{parskip}
\usepackage{wrapfig}
\usepackage[font=small,labelfont=bf]{caption}
% counting words? https://app.uio.no/ifi/texcount/online.php

\def\labelitemi{--}
\parindent=0pt

\begin{document}
\renewcommand{\refname}{\CHead{}}
% \setlength{\parindent}{0cm}
% \setlength{\parskip}{5pt}


{\bf Title of project:} Strengthening forest resilience through cross-disciplinary modeling of tree growth and decline across North America 

{\bf Starting date and ending date of the proposed fellowship:} 15 September 2027 to 29 February 2028 (24 weeks)

{\bf What will the visiting applicant bring to the host institution?} Wolkovich is a leader in the field of climate change ecology, with a cross-disciplinary perspective spanning ecology, climatology and evolutionary biology. She will thus help to strengthen links across these disciplines,  engaging with the departments of E3B (Ecology, evolution and environmental biology) and Earth and Environmental Sciences, while further bridging across the departments of Statistics and Political Science, where Gelman is appointed. She also plans to engage with the Columbia Climate School and researchers at Lamont-Doherty Earth Observatory bridging across campuses.

{\bf What will the host institution bring to the applicant?} Opportunity to bring fresh perspectives and new methods to modeling tree growth and forest health with a depth of expertise from statistics and workflow approaches, alongside best practices to communicate results for policy makers and the public. Gelman is a world leader in these topics and thus will be critical to developing novel modeling and interpretation methods. Additionally, the university's vibrant research community including the Columbia Climate School and Lamont-Doherty Earth Observatory will improve the design and impact of the research.

{\bf Description of project, objectives and methods:} Hierarchical modeling of tree growth across North American forests to improve our understanding of forest change over the last century and to the end of this century. This project will draw on data from radial tree growth data International Tree Ring Data Bank (ITRDB) alongside new fine-scale climate data to understand how individual-level tree growth scales up to lead to forest growth and decline, with implications for modeling temperate forests across the globe. As these forests can offset up to half of current fossil fuel emissions, understanding their responses is critical to resilience across ecosystems across the globe. 

\emph{Consider adding ...} They want to prioritize `new' collaborations (`new' not defined), so it would be good to say we have never worked intensively on a project before or such. 

\end{document}
